\section{Proof of the edge-contact classification}
\label{app:three:edge-classification}

Throughout this appendix, $q(x)=q_0+a^\top x+x^\top Qx$ satisfies the
hypotheses of \Cref{three:edge-classification}, and $d_i=\sqrt{Q_{ii}}>0$.
We assume that $q$ is not an affine square. Coordinate complementation
preserves vertex positivity, square coefficients, and extremality. We first
exclude a nonpositive mixed-coefficient product, then make every mixed
coefficient positive.

\subsection{The nonpositive-product sign class}

If $c_{12}c_{13}c_{23}\le0$, coordinate complementation makes all mixed
coefficients nonpositive, as in the proof of \Cref{three:sign}. The
exactness theorem of \citet[Theorem~1]{BurerNatarajanWillemsen2025} applies
to the SDP
\[
 \min\{\mathcal L_{m,Y}(q):
     \begin{pmatrix}1&m^\top\\m&Y\end{pmatrix}\succeq0,
     \ Y_{ij}\le m_i\ (i,j=1,2,3)\}.
\]
Its feasible set is compact. Indeed, the moment-matrix constraint gives
$Y_{ii}\ge m_i^2$, and $Y_{ii}\le m_i$ gives $0\le m_i\le1$.
The principal $2\times2$ moment blocks bound every
$|Y_{ij}|\le\sqrt{Y_{ii}Y_{jj}}\le1$. The uniform cube moments
$m_i=1/2$, $Y_{ii}=1/3$, $Y_{ij}=1/4$ for $i\ne j$ satisfy all inequalities
strictly and have a positive-definite moment matrix. Slater's condition
therefore gives an attained dual with the same value.

Put $\bar x=(1,x^\top)^\top$. Matching coefficients in the SDP
Lagrangian gives the dual
\[
 \max\left\{\beta:
 q(x)=\beta+\bar x^\top Z\bar x
             +\sum_{i,j=1}^3\lambda_{ij}x_i(1-x_j),\quad
 Z\succeq0,\ \lambda_{ij}\ge0\right\}.
\]
The displayed identity is a polynomial identity. Its optimum is
$\beta=\min_{C_3}q\ge0$ by the exactness theorem. Thus $q$ is a sum of
affine squares, a nonnegative constant, and nonnegative multiples of
$x_i(1-x_j)$, including the caps when $i=j$.

Extremality makes every nonzero summand proportional to $q$. Every product
$x_i(1-x_j)$ vanishes at a vertex, whereas $q$ is positive at all vertices,
so its multiplier must be zero. A nonzero constant cannot be proportional
to $q$ because $q$ has positive square coefficients. The PSD term therefore
makes $q$ an affine square, contrary to our assumption. Consequently the
mixed-coefficient product is positive, and a coordinate complementation
makes all three mixed coefficients positive for the rest of the proof.

\subsection{Zeros and perturbations}

A zero in the cube interior requires the full Hessian to be positive
semidefinite and the gradient to vanish. The quadratic is then a sum of
affine squares; extremality would make it a single affine square. A zero
in a facet interior is stationary in both facet coordinates and requires
the facet block to be positive semidefinite. This is excluded by
\Cref{three:stationary-facet}, proved in
\Cref{app:three:facet-classification}. Vertex zeros are excluded by
hypothesis. Thus every zero is in an edge interior, and there is at most
one on each edge because its quadratic coefficient is positive.

At a zero in an edge parallel to coordinate $i$, the derivative in that
coordinate vanishes. The two derivatives directed into adjacent facets
are nonnegative. They are in fact positive. If the inward derivative in
coordinate $j$ were zero, that zero would be stationary in both coordinates
of the adjacent facet. Its facet block would have to be positive
semidefinite: any negative quadratic direction can be reversed to point
inward in coordinate $j$, while coordinate $i$ admits both signs near the
edge-interior zero. \Cref{three:stationary-facet} excludes this stationary
facet zero for a non-square extreme ray.

Let $m$ be the number of edge zeros. Impose on a quadratic perturbation
$p$ its value and edge derivative being zero at each zero of $q$. These
are $2m$ homogeneous linear equations. Every such $p$ satisfies
$q\pm\epsilon p\ge0$ for sufficiently small $\epsilon>0$. To see this
locally, use an edge displacement $u$ and inward displacements $s,t\ge0$.
The positive edge curvature and strict inward derivatives give, in a
sufficiently small neighborhood,
\[
 q\ge \delta(u^2+s+t),\qquad
 |p|\le M(u^2+s+t)
\]
for some $\delta,M>0$. The second bound uses the zero value and edge
derivative of $p$; its other linear terms are bounded by multiples of
$s+t$, and its remaining quadratic terms by multiples of $u^2+s+t$.
The finitely many neighborhoods give a common permissible $\epsilon$.
On their compact complement, $q$ has a positive minimum, so reducing
$\epsilon$ also controls $p$ there.

If $m\le4$, the solution space of the perturbation equations has dimension
at least $10-2m\ge2$. It contains $q$ and a nonproportional $p$, contradicting
extremality. Consequently $m\ge5$. This argument uses an upper bound on
the number of independent equations; it does not assume that a contact
system has full rank.

\subsection{Two local restrictions on the contact graph}

For a vertex $v$, put $f_v=\sqrt{q(v)}>0$. If the edge $vw$ is parallel to
coordinate $i$ and has an interior zero, its restriction is a square and
\begin{equation}\label{eq:three:edge-length}
 f_v+f_w=d_i.
\end{equation}
Indeed its zero divides the unit edge into lengths $f_v/d_i$ and
$f_w/d_i$. Conversely, \eqref{eq:three:edge-length} with both endpoint
values positive implies an interior zero: the endpoint values and the
quadratic coefficient determine that square restriction.

Suppose two contact edges at $v$ use coordinates $i,j$. Reflect their
coordinates so that both point positively from $v$. Writing their zero
positions as $\alpha_i=f_v/d_i$ and $\alpha_j=f_v/d_j$, the restriction to
their face has the form
\[
 f_v^2\left(1-\frac{u_i}{\alpha_i}
                 -\frac{u_j}{\alpha_j}\right)^2
  + (\widehat c_{ij}-2d_id_j)u_iu_j,
\]
where $\widehat c_{ij}=(-1)^{v_i+v_j}c_{ij}$.
At $(\alpha_i/2,\alpha_j/2)$, nonnegativity implies
$\widehat c_{ij}\ge2d_id_j$. Since $c_{ij}>0$, this forces $v_i=v_j$.
The inequality need not be strict at this stage.

Three contact edges at one vertex are also impossible for an extreme
$q$. The same calculation in three coordinates gives an identity
\[
 q=f_v^2\left(1-\sum_i\frac{u_i}{\alpha_i}\right)^2
       +\sum_{i<j}(\widehat c_{ij}-2d_id_j)u_iu_j.
\]
Every summand is nonnegative on the reflected cube. If every product
coefficient is zero, $q$ is the excluded affine square. Otherwise this is
a nontrivial decomposition of $q$ into nonnegative quadratics, contrary
to extremality.

We shall also use a two-variable observation. Suppose a nonnegative
quadratic on a square has positive mixed coefficient and has interior
zeros on two opposite edges parallel to coordinate $i$. Let their zero
positions be $\alpha,\beta\in(0,1)$, at coordinate $j=0,1$, respectively.
Then
\[
 q(u_i,u_j)=d_i^2\bigl(u_i-\alpha-(\beta-\alpha)u_j\bigr)^2
                 +\eta u_j(1-u_j),\qquad
 \eta=d_i^2(\beta-\alpha)^2-d_j^2.
\]
The mixed coefficient being positive gives $\alpha>\beta$. Evaluating at
the affine minimizer in $u_i$, which stays in $(0,1)$, gives $\eta\ge0$.
Thus
\begin{equation}\label{eq:three:opposite-order}
 d_j\le d_i(\alpha-\beta)<d_i.
\end{equation}

\subsection{The possible contact graphs}

Let the Hamming weight of a vertex be its number of unit coordinates.
The rule $v_i=v_j$ shows that, when two contact edges meet, their other
endpoints have equal Hamming weight. Hence a connected component of the
contact graph stays between one pair of adjacent weight levels. The cube
edges between levels $0,1$ form a three-edge star, those between levels
$1,2$ form a six-edge cycle, and those between levels $2,3$ form a second
three-edge star. Every vertex has contact degree at most two.

Write $a,b$ for the numbers of contact edges in the lower and upper stars,
and $m_0$ for the number in the central cycle. Thus $0\le a,b\le2$ and
$a+b+m_0\ge5$. A leaf of an outer star cannot also meet a central contact
edge, because that would join different pairs of weight levels. Up to
coordinate permutation and simultaneous complementation of all three
coordinates, the following cases exhaust the possibilities.

\begin{enumerate}[label=(\roman*)]
\item $a=b=0$. There are at least five central-cycle edges.
\item $a=1,b=0$. Deleting the level-one endpoint of the lower edge leaves
      a four-edge path in the central cycle. All four edges must be present.
\item $a=b=1$. Deleting their central endpoints leaves at most a three-edge
      path. To reach five contacts, the endpoints must be adjacent in the
      cycle and all three remaining edges must be present.
\item $a=2,b=1$. Only one level-one vertex remains available centrally, so
      the central contacts must be its two incident edges. The upper-star
      endpoint is its nonneighbor in level two.
\item $a=b=2$. Only one vertex in each central level remains available.
      They must be adjacent, and their edge must be present.
\end{enumerate}
The case $a=2,b=0$ cannot reach five contacts, because at most two central
edges remain available. Interchanging the outer stars accounts for the
cases with $b>a$ not listed above. These observations also cover graphs
with more than five contacts: the central-only case allows six, and the
availability bounds in every other case allow at most five.

\subsection{Excluding the central cycle and three other cases}

In case (i), five equations \eqref{eq:three:edge-length} force the sixth.
To verify this, list the six directions around the central cycle as
$1,3,2,1,3,2$. The alternating sum of the six left sides is zero, and so
is the alternating sum of their right sides. Both endpoint square roots
are positive, so the missing edge also has an interior zero.

Set $h=f_{100}+d_1$. Applying \eqref{eq:three:edge-length} at the shared
level-two vertices gives
\[
 f_{e_i}=h-d_i,\qquad
 f_{e_i+e_j}=d_i+d_j-h.
\]
It follows that $q$ and $(d_1x_1+d_2x_2+d_3x_3-h)^2$ have the same square
coefficients and the same values at the six central vertices. Their
difference is therefore
\[
 A r(x),\qquad
 r(x)=1-x_1-x_2-x_3+x_1x_2+x_1x_3+x_2x_3.
\]
For completeness, a multilinear quadratic vanishing at $e_i$ has linear
coefficients equal to the negative of its constant; vanishing at
$e_i+e_j$ then makes every mixed coefficient equal to that constant.
At a central edge zero, the square has zero gradient and $r$ has strictly
positive inward derivatives. Thus nonnegativity forces $A\ge0$.
Since
\[
 r=(1-x_1)(1-x_2)(1-x_3)+x_1x_2x_3,
\]
this decomposes $q$ into an affine square and an element of $\mathcal D_3$.
Extremality would force $q$ to be the excluded affine square.

For case (iii), choose the lower edge as $000$--$100$ and the upper edge
as $110$--$111$. The three central edges are
\[
 001\text{--}101,\qquad001\text{--}011,\qquad010\text{--}011.
\]
The face $x_2=0$ has opposite contact edges parallel to $x_1$, so
\eqref{eq:three:opposite-order} gives $d_3<d_1$. The face $x_2=1$ has
opposite contact edges parallel to $x_3$, giving $d_1<d_3$, a contradiction.

For case (iv), take the lower edges as $000$--$100$ and $000$--$010$.
The two central edges are $001$--$101$ and $001$--$011$, and the upper
edge is $110$--$111$. The two adjacent contacts on $x_3=0$ give
\[
 q(x_1,x_2,0)=(h-d_1x_1-d_2x_2)^2+T x_1x_2,
 \quad 0<h<\min(d_1,d_2),\quad T\ge0.
\]
The corresponding contacts on $x_3=1$ give the same formula with $h$
replaced by $g$, where $0<g<\min(d_1,d_2)$. The mixed coefficient is
unchanged, so the same $T$ appears. Matching the square coefficient of
$x_3$ now gives
\[
 q=\bigl(h+(g-h)x_3-d_1x_1-d_2x_2\bigr)^2
          +T x_1x_2+\eta x_3(1-x_3),\qquad
 \eta=(g-h)^2-d_3^2\ge0.
\]
The last inequality follows from the two opposite $x_1$ contacts on
$x_2=0$. On the proposed upper contact edge $x_1=x_2=1$, the square is
strictly positive, because the affine interpolation of $h,g$ is less
than $d_1+d_2$. The other terms are nonnegative, a contradiction.

For case (v), choose the lower edges as $000$--$100$ and $000$--$010$,
the upper edges as $011$--$111$ and $110$--$111$, and the central edge
as $001$--$101$. There are three contacts parallel to $x_1$, at
$(x_2,x_3)=(0,0),(0,1),(1,1)$. Complete the square in $x_1$:
\[
 q=d_1^2\bigl(x_1-\alpha(x_2,x_3)\bigr)^2+R(x_2,x_3),\qquad
 \alpha(x_2,x_3)=\alpha_0-u x_2-v x_3,
\]
where $u,v>0$ because the two mixed coefficients involving $x_1$ are
positive. Each of the three recorded values of $\alpha$ lies in $(0,1)$.
Its fourth vertex value lies between $\alpha(1,1)$ and $\alpha(0,0)$,
so $\alpha$ lies in $(0,1)$ throughout the square. Consequently $R\ge0$
there. On the lower contact edge $x_1=x_3=0$, however,
$\alpha(x_2,0)>0$, so the displayed square is strictly positive and
$q$ cannot vanish. This is again a contradiction.

\subsection{Recovering the family from the remaining graph}

Case (ii) is the only possibility left. For the representative with isolated
edge $000$--$100$, its central path is
$110$--$010$--$011$--$001$--$101$. Swap coordinates $1$ and $3$, then
complement coordinates $1$ and $2$. This puts the contact edges in the
following positions:
\[
 000\text{--}100,\quad000\text{--}010,\quad
 100\text{--}101,\quad010\text{--}011,\quad110\text{--}111.
\]
The transformed quadratic need no longer have all mixed coefficients
positive, but its square coefficients remain $d_i^2$.
Let $h=\sqrt{q(000)}>0$. The first two contacts fix the edge roots
$h/d_1,h/d_2\in(0,1)$, the constant $h^2$, and the linear coefficients
$-2hd_1,-2hd_2$. The next two contacts fix the roots
$(d_1-h)/d_3$ and $(d_2-h)/d_3$.
Write the fifth root as $r/d_3\in(0,1)$, and set
$D=d_1+d_2-h$ and $k=r-D$.

The derivative equations on the three vertical edges and the value on
the edge through $110$ determine the remaining coefficients:
\[
 a_3=2d_3(h+k),\qquad
 c_{13}=-2d_3(d_1+k),\qquad
 c_{23}=-2d_3(d_2+k),\qquad
 c_{12}=2d_1d_2+k(2D+k).
\]
Thus $q=q_{h,d,k}$. At $(h/d_1,0,0)$ the inward derivative in $x_3$ is
$2d_3k(1-h/d_1)$. This derivative is strictly positive by the first part
of the proof, so $k>0$. The fifth root gives $D+k<d_3$,
while the first two give $0<h<\min(d_1,d_2)$. These are exactly the strict
five-contact assumptions.

Conversely, such a family member generates an exposed extreme ray by the
five-contact argument. Its vertices have positive values: the contact
edge endpoints are positive, and the remaining vertex $001$ has value
$(h+d_3)^2+2d_3k>0$. Its mixed coefficient product is positive and each
$2\times2$ principal determinant is negative, as follows from
$D>0$, $k>0$, and the coefficient formulas above.

\subsection{The affine-square rays}

Finally let $q=\ell^2$ have positive values at every cube vertex. If the
affine factor $\ell$ takes only one sign there, it is bounded away from
zero on the cube. Then $q$ is strictly positive and lies in the interior
of $\mathcal P_3$, so its ray is not extreme. If $\ell$ takes both signs,
its zero plane has a relatively open patch in the cube interior.
In any decomposition $q=p+r$ with $p,r\in\mathcal P_3$, both summands
vanish on this patch. Polynomial divisibility gives $p=\ell\ell_p$ for
an affine $\ell_p$. Nonnegativity on both sides of the plane forces
$\ell_p$ to vanish on the patch; otherwise the product changes sign
near a point where $\ell_p\ne0$. Hence $\ell_p=\alpha\ell$, and
evaluating away from the plane gives $\alpha\ge0$. The same argument
applies to $r$. Thus the ray of $\ell^2$ is extreme. This completes
the proof of \Cref{three:edge-classification}.
